\section{Architecture du projet}

L'architecture de \texttt{Quoridor} repose sur une organisation en couches.
Ce choix est adapté à la nature du logiciel : le projet ne se limite pas à
une interface unique, mais combine un CLI interactif, un mode contest, une
GUI GTK et un mode réseau client/serveur. Dans ce contexte, la priorité
architecturale n'est pas de structurer l'application autour d'un paradigme
graphique particulier, mais de centraliser la logique métier afin qu'elle soit
réutilisable depuis plusieurs points d'entrée.

L'objectif principal de cette architecture est double. D'une part, il faut
isoler le coeur du jeu afin que les règles, les structures de données et les
algorithmes puissent être testés indépendamment des interfaces. D'autre part,
il faut permettre à plusieurs modes d'exécution de partager les mêmes services
applicatifs sans dupliquer la logique de validation, de persistance ou
d'intelligence artificielle. Cette organisation répond donc à une exigence de
cohérence fonctionnelle autant qu'à une exigence de maintenabilité.

\subsection{Vue d'ensemble}

Le point d'entrée du programme est le module \texttt{quoridor.\_\_main\_\_},
qui délègue l'exécution au CLI principal. Celui-ci interprète les options,
charge la configuration, initialise l'internationalisation, puis redirige
l'exécution vers le bon sous-système. L'architecture peut être résumée en six
blocs principaux :
\begin{itemize}
  \item les \textbf{interfaces}, qui collectent les actions utilisateur et
    affichent les résultats ;
  \item la \textbf{couche applicative}, qui orchestre les cas d'usage ;
  \item le \textbf{noyau métier}, qui contient l'état du jeu et les objets de
    domaine ;
  \item la \textbf{couche règles}, qui encode les contraintes de Quoridor ;
  \item les \textbf{services techniques transverses}, comme la configuration,
    l'internationalisation et la persistance ;
  \item l'\textbf{infrastructure réseau et IA}, qui étend l'application sans
    modifier le coeur du jeu.
\end{itemize}

La Fig.~\ref{fig:architecture-layered-overview} présente cette organisation
générale.

\begin{figure}[!ht]
  \centering
  \begin{tikzpicture}[
      >=Latex,
      font=\small,
      package/.style={
        draw,
        rounded corners,
        align=left,
        minimum width=4.1cm,
        minimum height=2.1cm,
        fill=blue!5
      },
      support/.style={
        draw,
        rounded corners,
        align=left,
        minimum width=3.5cm,
        minimum height=1.6cm,
        fill=gray!10
      },
      arrow/.style={->, thick}
    ]
    \node[package] (interfaces) {
      \textbf{Interfaces}\\
      \texttt{cli.py}\\
      \texttt{cli\_shell.py}\\
      \texttt{contest\_parser.py}\\
      \texttt{gui.py}\\
      \texttt{cli\_network.py}
    };

    \node[package, below=1.2cm of interfaces] (application) {
      \textbf{Application}\\
      \texttt{game\_application\_service.py}\\
      \texttt{game\_session.py}\\
      \texttt{history\_manager.py}\\
      \texttt{blitz.py}\\
      \texttt{persistence\_service.py}
    };

    \node[package, below left=1.2cm and 0.5cm of application] (core) {
      \textbf{Core}\\
      \texttt{game\_state.py}\\
      \texttt{board.py}\\
      \texttt{notation.py}\\
      \texttt{move\_record.py}
    };

    \node[package, below=1.2cm of application] (rules) {
      \textbf{Rules}\\
      \texttt{pawn\_rules.py}\\
      \texttt{wall\_rules.py}\\
      \texttt{win\_rules.py}\\
      \texttt{pathfinding.py}
    };

    \node[package, below right=1.2cm and 0.5cm of application] (utils) {
      \textbf{Utils}\\
      \texttt{graph.py}
    };

    \node[support, right=1.4cm of application] (network) {
      \textbf{Réseau}\\
      \texttt{network/server.py}\\
      \texttt{network/client.py}\\
      \texttt{network/discovery.py}
    };

    \node[support, left=1.4cm of application] (crosscutting) {
      \textbf{Services transverses}\\
      \texttt{config.py}\\
      \texttt{i18n.py}
    };

    \node[support, right=1.4cm of rules] (ml) {
      \textbf{IA avancée / ML}\\
      \texttt{minimax\_engine.py}\\
      \texttt{mcts\_engine.py}\\
      \texttt{ML/*}
    };

    \draw[arrow] (interfaces) -- (application);
    \draw[arrow] (application) -- (core);
    \draw[arrow] (application) -- (rules);
    \draw[arrow] (rules) -- (utils);
    \draw[arrow] (network) -- (application);
    \draw[arrow] (ml) -- (application);
    \draw[arrow] (crosscutting) -- (interfaces);
    \draw[arrow] (crosscutting) -- (application);
  \end{tikzpicture}
  \caption{Diagramme UML de packages simplifié de l'architecture de Quoridor}
  \label{fig:architecture-layered-overview}
\end{figure}

\subsection{Découpage modulaire}

\paragraph{Couche \texttt{interfaces}.}
Cette couche regroupe tous les points d'entrée visibles par l'utilisateur. Le
module \texttt{cli.py} sert de routeur principal entre le mode interactif, le
mode contest, le mode serveur et la GUI. Le module \texttt{cli\_shell.py}
porte la boucle interactive historique, mais il a été progressivement
factorisé à l'aide du sous-paquet \texttt{interfaces.shell}, qui regroupe des
responsabilités spécialisées : amorçage de session, état courant du shell,
gestion du blitz, commandes locales, commandes réseau et événements. La GUI
est concentrée dans \texttt{gui.py}, qui partage la logique métier avec le
CLI via les services applicatifs plutôt que de réimplémenter les règles.

\paragraph{Couche \texttt{application}.}
Cette couche orchestre les cas d'usage complets. La classe
\texttt{GameSession} encapsule la dynamique d'une partie : application d'un
coup, alternance des joueurs, calcul de vainqueur, historique, undo/redo et
intégration du blitz. La classe \texttt{GameApplicationService} joue le rôle
de façade : elle expose aux interfaces des opérations de haut niveau comme
charger, sauvegarder, demander un conseil, jouer un coup au format texte ou
créer une nouvelle session. Cette séparation limite fortement le couplage
entre l'interface et le coeur du jeu.

\paragraph{Couche \texttt{core}.}
Le noyau métier repose principalement sur \texttt{GameState}, qui constitue la
source de vérité unique pour l'état courant d'une partie. Cette structure
regroupe la taille du plateau, le joueur courant, les positions des pions, les
murs restants, les murs posés, ainsi que les joueurs temporairement inactifs
en mode réseau. Cette couche contient également la notation, les snapshots et
les objets nécessaires à la sérialisation de l'historique.

\paragraph{Couche \texttt{rules}.}
Les règles du jeu sont isolées dans des modules dédiés : déplacement des
pions, validation des murs, victoire et recherche de chemin. Cette couche est
la plus proche du domaine pur : elle est largement stateless et déterministe.
Elle est donc particulièrement testable et réutilisable. Elle manipule
directement la structure de graphe sous-jacente, ce qui fait le lien naturel
entre architecture et structures de données.

\paragraph{Couche \texttt{utils}.}
Le module \texttt{utils.graph} fournit la structure de graphe utilisée pour
représenter la connectivité du plateau. Le plateau est modélisé comme une
liste d'adjacence dont les arêtes sont supprimées lors de la pose des murs.
Cette structure technique est volontairement maintenue en dehors des règles
elles-mêmes, afin de rendre son implémentation réutilisable.

\paragraph{Infrastructure réseau.}
Le sous-paquet \texttt{network} implémente un client, un serveur multi-clients
et un mécanisme de découverte UDP des serveurs locaux. D'un point de vue
architectural, il ne constitue pas un second moteur de jeu : il s'appuie sur
la même logique métier que le mode local. Le serveur utilise lui aussi les
objets \texttt{GameSession} et les services applicatifs pour valider les coups
et maintenir les parties.

\paragraph{Infrastructure IA et ML.}
Les moteurs \texttt{minimax} et \texttt{mcts} appartiennent à la couche
applicative car ils exploitent le modèle métier pour produire des coups. Le
sous-dossier \texttt{ML} contient les artefacts liés à l'entraînement et à
l'extraction de caractéristiques pour le sélecteur ML utilisé avec MCTS.
Même si cette partie reste moins mature que le reste du projet, elle est bien
identifiée comme une extension algorithmique du coeur du jeu.

\subsection{Dépendances entre modules}

La direction générale des dépendances est descendante. Les interfaces
dépendent de la couche applicative, la couche applicative dépend du noyau
métier et des règles, et les règles s'appuient elles-mêmes sur les structures
de données techniques. Cette orientation est saine car elle évite qu'une règle
du jeu ne dépende d'un choix d'affichage, d'une entrée clavier ou d'une
fenêtre graphique.

Le module \texttt{GameApplicationService} est le principal point de couture de
cette architecture. Il masque la complexité interne du moteur à l'interface
appelante. Ainsi, la GUI et le CLI invoquent les mêmes méthodes pour charger
une partie, sauvegarder, jouer un coup ou demander un conseil. Cette façade
limite la dispersion de la logique métier et facilite l'introduction de tests
fonctionnels communs.

Certaines dépendances méritent cependant d'être discutées de manière critique.
Tout d'abord, le shell interactif historique reste un composant volumineux,
même après factorisation. Ensuite, la persistance dépend encore d'un parseur
de mode contest situé dans la couche \texttt{interfaces}, ce qui n'est pas
idéal d'un point de vue purement architectural. Enfin, le moteur MCTS importe
directement le sélecteur ML au runtime, ce qui crée un couplage plus fort que
nécessaire entre l'algorithme standard et son extension d'apprentissage.

\subsection{Principales interfaces internes}

Plutôt que de décrire exhaustivement toutes les fonctions, il est plus utile
de mettre en avant les interfaces internes qui structurent réellement le
projet.

\begin{itemize}
  \item \textbf{\texttt{GameSession}} : interface centrale pour appliquer les
    coups, changer de tour, manipuler l'historique et calculer le gagnant.
    Cette classe expose notamment les opérations
    \texttt{play\_pawn\_move\_from\_to()}, \texttt{place\_wall()},
    \texttt{undo()}, \texttt{redo()} et \texttt{compute\_ai\_move()}.
  \item \textbf{\texttt{GameApplicationService}} : façade applicative exposée
    au CLI et à la GUI. Elle fournit des points d'entrée cohérents comme
    \texttt{load()}, \texttt{save()}, \texttt{hint()},
    \texttt{play\_pawn\_move\_token()} et \texttt{place\_wall\_token()}.
  \item \textbf{\texttt{NetworkServer} / \texttt{NetworkClient}} : interface du
    mode réseau, responsable du protocole et du transport, mais pas des règles
    métier elles-mêmes.
  \item \textbf{\texttt{GameState}} : représentation canonique de l'état du
    jeu, utilisée par toutes les couches internes.
  \item \textbf{\texttt{ConfigManager} et \texttt{setup\_i18n}} : services
    transverses partagés par plusieurs interfaces.
\end{itemize}

Ce choix d'interfaces est important d'un point de vue architectural. Les
interfaces utilisateur ne manipulent pas directement les structures profondes
du domaine ; elles passent par des services ou des sessions explicitement
identifiés. En pratique, cela signifie qu'une commande du shell, un clic dans
la GUI ou un événement réseau peuvent déclencher des opérations différentes au
niveau de l'interface, tout en convergeant vers les mêmes abstractions
internes. Cette convergence améliore la cohérence fonctionnelle et réduit le
risque de divergence entre les modes d'utilisation du logiciel.

La Fig.~\ref{fig:architecture-key-components} illustre ces composants et leurs
interactions majeures.

\begin{figure}[!ht]
  \centering
  \begin{tikzpicture}[
      >=Latex,
      font=\small,
      class/.style={
        draw,
        rounded corners,
        align=left,
        minimum width=4.0cm,
        minimum height=1.5cm,
        fill=green!6
      },
      infra/.style={
        draw,
        rounded corners,
        align=left,
        minimum width=3.6cm,
        minimum height=1.3cm,
        fill=orange!10
      },
      arrow/.style={->, thick}
    ]
    \node[class] (cli) {
      \textbf{CLI / Shell}\\
      \texttt{cli.py}\\
      \texttt{cli\_shell.py}\\
      \texttt{interfaces.shell.*}
    };

    \node[class, right=1.5cm of cli] (gui) {
      \textbf{GUI GTK}\\
      \texttt{gui.py}\\
      \texttt{gui\_shortcuts.py}
    };

    \path (cli) -- (gui) coordinate[midway] (topmid);

    \node[class, below=1.4cm of topmid] (service) {
      \textbf{GameApplicationService}\\
      chargement, sauvegarde, hint,\\
      validation des commandes métier
    };

    \node[class, below left=1.6cm and 0.6cm of service] (session) {
      \textbf{GameSession}\\
      tours, historique, undo/redo,\\
      intégration du blitz
    };

    \node[class, below right=1.6cm and 0.6cm of service] (state) {
      \textbf{GameState}\\
      positions, murs, joueur courant,\\
      snapshots
    };

    \node[infra, left=1.4cm of service] (network) {
      \textbf{Réseau}\\
      \texttt{NetworkClient}\\
      \texttt{NetworkServer}
    };

    \node[infra, right=1.4cm of service] (ai) {
      \textbf{IA}\\
      minimax, iterative, MCTS,\\
      sélecteur ML
    };

    \node[infra, below=1.4cm of state] (rules) {
      \textbf{Rules + Graph}\\
      validation des coups,\\
      accessibilité, victoire
    };

    \draw[arrow] (cli) -- (service);
    \draw[arrow] (gui) -- (service);
    \draw[arrow] (network) -- (service);
    \draw[arrow] (service) -- (session);
    \draw[arrow] (service) -- (ai);
    \draw[arrow] (session) -- (state);
    \draw[arrow] (session) -- (rules);
    \draw[arrow] (state) -- (rules);
  \end{tikzpicture}
  \caption{Diagramme UML de composants simplifié et dépendances principales}
  \label{fig:architecture-key-components}
\end{figure}

\subsection{Flux de contrôle typique}

Pour rendre l'architecture plus concrète, on peut suivre le traitement d'une
commande simple du CLI, par exemple un déplacement de pion. Le shell
interactif lit la commande utilisateur, la normalise, puis délègue son
interprétation à ses handlers spécialisés. Ceux-ci convertissent la saisie en
appel vers la couche applicative ou directement vers \texttt{GameSession} selon
le cas. La session consulte alors \texttt{GameState}, sollicite les modules de
règles pour valider le coup, met à jour l'état, enregistre l'action dans
l'historique, puis rend le contrôle à l'interface qui réaffiche l'état du jeu.

Le même principe vaut pour la GUI et pour le réseau. Dans la GUI, l'action
provient d'un menu ou d'un glisser-déposer ; côté réseau, elle provient d'une
commande reçue par le serveur. Dans les deux cas, la validation du coup repose
sur les mêmes modules métier. Ce point est essentiel : l'architecture ne se
contente pas de juxtaposer plusieurs interfaces, elle leur impose un noyau
commun qui garantit un comportement homogène.

\subsection{Lien avec les structures de données et les algorithmes}

L'architecture choisie n'est pas indépendante des structures de données.
Au contraire, elle a été construite pour leur laisser une place centrale sans
les exposer directement aux interfaces. La structure de graphe utilisée pour
représenter le plateau est manipulée par \texttt{GameState} et exploitée par
les modules de règles. Les algorithmes de recherche de chemin, de validation
des murs et de calcul des coups légaux sont donc situés là où leur utilité est
la plus naturelle : au voisinage immédiat du domaine.

Les algorithmes d'intelligence artificielle utilisent ensuite ces mêmes
structures de données via des services applicatifs. Minimax, iterative
deepening et MCTS ne sont pas branchés directement sur une interface, mais sur
l'état du jeu et les règles. Ce positionnement garantit qu'un conseil affiché
en CLI, une suggestion de coup en GUI ou une décision prise en mode réseau
s'appuient sur la même base algorithmique.

Ce lien explicite entre architecture et structures de données justifie le
découpage retenu. Une architecture uniquement pensée autour des interfaces
aurait eu tendance à disperser la logique de validation dans plusieurs
composants, alors que le choix effectué place la représentation du plateau et
les algorithmes de graphe au centre du système. C'est précisément ce point qui
rend possible la réutilisation des mêmes règles pour le jeu local, l'undo/redo,
la sauvegarde, l'IA et le mode réseau.

\subsection{Justification des choix}

Les choix architecturaux retenus répondent principalement à trois objectifs.
Le premier est la \textbf{réutilisation} : une même logique métier doit
alimenter plusieurs interfaces sans divergence de comportement. Le second est
la \textbf{testabilité} : il doit être possible de tester le coeur du projet
sans démarrer l'interface graphique ni le réseau. Le troisième est
l'\textbf{extensibilité} : l'ajout du blitz, du réseau ou de nouvelles
heuristiques IA doit rester possible sans refondre complètement le coeur du
jeu.

Une autre justification importante tient à la granularité choisie. Le projet
aurait pu rester organisé autour de quelques fichiers très larges, notamment
pour le shell ou les services applicatifs. Au contraire, la factorisation en
sous-modules spécialisés permet de limiter la taille des composants, de
clarifier les responsabilités et de rendre les tests plus ciblés. Cette
approche n'élimine pas complètement la complexité, mais elle la rend plus
maîtrisable.

Cette architecture présente toutefois certaines limites. Le shell interactif
reste historiquement central et a nécessité une factorisation progressive.
L'intégration GUI n'offre pas encore la même surface fonctionnelle que le CLI,
notamment pour le réseau. Enfin, la partie ML n'est pas encore aussi bien
isolée que les autres briques. Ces points n'annulent pas la cohérence
d'ensemble, mais ils constituent des axes d'amélioration réalistes.

\subsection{Bilan architectural}

Au final, l'architecture de \texttt{Quoridor} est lisible, implémentable et
cohérente avec la complexité réelle du projet. Elle évite de placer les règles
du jeu dans les interfaces, centralise l'état dans une structure métier
identifiée, et s'appuie sur des services applicatifs communs pour partager les
cas d'usage entre le CLI, la GUI et le mode réseau. C'est cette combinaison
entre séparation des responsabilités, réutilisation et lien explicite avec les
algorithmes du jeu qui en fait une base solide pour le projet final.
